\documentclass{article}
\usepackage[utf8]{inputenc}
\usepackage[T1]{fontenc}
\usepackage[indonesian]{babel}
\usepackage{hyperref}
\usepackage[a4paper,margin=2.5cm]{geometry}
\usepackage{booktabs}
\usepackage{tabularx}
\usepackage{array}
\usepackage{float}
\usepackage{tikz}
\usepackage{amsmath}
\usetikzlibrary{arrows.meta,positioning,shapes.geometric}
\usepackage{xcolor}
\usepackage{listings}

\lstdefinelanguage{Rust}{
  morekeywords={as,break,const,continue,crate,else,enum,extern,false,fn,for,if,impl,in,let,loop,match,mod,move,mut,pub,ref,return,self,Self,static,struct,super,trait,true,type,unsafe,use,where,while,async,await,dyn},
  sensitive=true,
  morecomment=[l]{//},
  morecomment=[s]{/*}{*/},
  morestring=[b]",
}

\lstdefinestyle{ruststyle}{
  language=Rust,
  basicstyle=\ttfamily\small,
  keywordstyle=\color{blue!70!black}\bfseries,
  commentstyle=\color{green!40!black}\itshape,
  stringstyle=\color{orange!70!black},
  numbers=left,
  numberstyle=\tiny\color{gray},
  numbersep=8pt,
  breaklines=true,
  frame=single,
  showstringspaces=false,
  tabsize=2,
  captionpos=b
}

\title{ALGORUST: Dokumen Teknis Pembelajaran Algoritma, P2P, dan Kriptografi Dasar Berbasis Rust}
\author{MagerCode Team \\ \small \texttt{contact@rayhex.biz.id}}
\date{\today}

\begin{document}
\maketitle

\section{Ringkasan Eksekutif}
ALGORUST adalah kerangka belajar teknis berbahasa Indonesia yang dirancang untuk menjembatani teori
algoritma dengan implementasi perangkat lunak nyata menggunakan Rust. Dokumen ini menjelaskan
arah pembelajaran, struktur modul, metodologi evaluasi, serta praktik keamanan minimum
yang perlu dipahami peserta.

\section{Sejarah dan Filosofi Bahasa Rust}
Bahasa pemrograman Rust pertama kali dikembangkan oleh Mozilla Research sebagai proyek pribadi
Graydon Hoare pada tahun 2006, dan secara resmi diumumkan oleh Mozilla pada tahun 2010.
Rust dirancang untuk menjadi bahasa sistem yang aman secara memori (\textit{memory-safe}) tanpa
mengorbankan kecepatan atau membutuhkan pengumpul sampah (\textit{garbage collector}).

Filosofi utama Rust adalah memberikan kontrol tingkat rendah seperti C/C++, namun dengan jaminan
keamanan memori melalui sistem kepemilikan (\textit{ownership}) dan pemeriksaan pinjaman (\textit{borrow checker}).
Sejak rilis stabil pertama pada tahun 2015, Rust tumbuh pesat dan dinobatkan sebagai bahasa
``paling dicintai'' dalam survei Stack Overflow selama bertahun-tahun. Saat ini Rust digunakan
secara luas di berbagai domain, termasuk pengembangan sistem operasi, mesin browser, perangkat
tertanam, dan tentu saja aplikasi jaringan terdesentralisasi seperti blockchain dan protokol P2P.

Keandalan dan kecepatan Rust menjadikannya pilihan tepat untuk mengimplementasikan algoritma
kriptografi dan protokol jaringan yang membutuhkan ketepatan tinggi dan performa maksimal.

\section{Visi dan Tujuan}
\begin{itemize}
  \item Menyediakan jalur belajar terstruktur dari tingkat dasar hingga menengah-lanjut.
  \item Menjelaskan keputusan desain algoritma berdasarkan kompleksitas waktu, ruang, dan keterbacaan kode.
  \item Membentuk kebiasaan rekayasa perangkat lunak yang terukur melalui pengujian, analisis, dan refleksi solusi.
  \item Menyediakan materi yang relevan untuk kelas rekayasa perangkat lunak maupun pembelajaran mandiri.
\end{itemize}

\section{Ruang Lingkup Pembelajaran}
\subsection{Modul 1: Algoritma Fundamental}
Modul ini membahas pengurutan, pencarian, penelusuran graf, dan pemrograman dinamis sebagai fondasi
berpikir komputasional yang efisien.

\begin{table}[H]
\centering
\caption{Ringkasan Topik Algoritma Fundamental}
\begin{tabularx}{\linewidth}{>{\raggedright\arraybackslash}p{3cm} >{\raggedright\arraybackslash}X >{\raggedright\arraybackslash}p{3cm}}
\toprule
Kategori & Materi & Kompleksitas Kunci \\
\midrule
Sorting & Bubble, Selection, Insertion, Merge, Quick Sort & Perbandingan $O(n^2)$ vs $O(n \log n)$ \\
Pencarian & Linear Search, Binary Search & Trade-off data terurut dan tidak terurut \\
Graf & BFS, DFS, Dijkstra & Penelusuran $O(V+E)$ dan lintasan terpendek \\
Pemrograman Dinamis & Knapsack, Coin Change & Optimasi submasalah dan penyimpanan hasil antara (memoization) \\
\bottomrule
\end{tabularx}
\end{table}

Rumus rekursif untuk kompleksitas Merge Sort adalah:
\[
T(n) = 2T\left(\frac{n}{2}\right) + O(n)
\]
yang menghasilkan solusi \(T(n) = O(n \log n)\).

\subsection{Modul 2: Algoritma Jaringan P2P}
% ============= TAMBAHAN: SEJARAH DAN KONSEP P2P =============
Jaringan \textit{Peer-to-Peer} (P2P) pertama kali dikenal luas melalui aplikasi berbagi berkas
seperti Napster (1999) dan BitTorrent (2001). Konsep utamanya adalah setiap node (rekan)
dapat berperan sebagai klien sekaligus server, sehingga tidak bergantung pada server pusat.
Hal ini memberikan ketahanan dan skalabilitas tinggi. Perkembangan terbaru teknologi blockchain
seperti Bitcoin dan Ethereum juga menggunakan prinsip P2P untuk mencapai konsensus terdesentralisasi.

Fokus utama modul ini adalah simulasi penyebaran pesan dan pengambilan keputusan antar node
secara terdesentralisasi.

\begin{table}[H]
\centering
\caption{Skenario Simulasi P2P pada ALGORUST}
\begin{tabularx}{\linewidth}{>{\raggedright\arraybackslash}p{3cm} >{\raggedright\arraybackslash}X >{\raggedright\arraybackslash}p{3cm}}
\toprule
Skenario & Deskripsi & Metrik Evaluasi \\
\midrule
Penyiaran Gosip & Tiap node meneruskan pesan ke beberapa rekan acak & Latensi propagasi \\
Pemeringkatan Rekan & Node memberi bobot reputasi pada rekan & Stabilitas jaringan \\
Fault Injection & Simulasi node gagal atau lambat & Ketahanan protokol \\
\bottomrule
\end{tabularx}
\end{table}

Dalam protokol gosip, probabilitas sebuah node menerima pesan setelah \(k\) putaran dapat dinyatakan dengan:
\[
P_{\text{terima}} \approx 1 - (1 - p)^{k \cdot d}
\]
dengan \(p\) adalah probabilitas forwarding dan \(d\) adalah derajat rata-rata node.

\subsection{Modul 3: Kriptografi Dasar}
Kriptografi telah digunakan sejak zaman Romawi dengan cipher sederhana seperti Caesar Cipher.
Perkembangan pesat terjadi pada abad ke-20, terutama saat Perang Dunia II dengan mesin Enigma.
Era modern ditandai dengan ditemukannya kriptografi kunci publik oleh Diffie dan Hellman (1976)
yang memungkinkan pertukaran kunci aman melalui kanal tidak aman. Saat ini kriptografi menjadi
fondasi keamanan digital, termasuk dalam mata uang kripto dan protokol komunikasi.

Modul ini menekankan pemahaman konsep inti keamanan informasi melalui implementasi sederhana
yang mudah dipelajari.

\begin{table}[H]
\centering
\caption{Modul Kriptografi Dasar}
\begin{tabularx}{\linewidth}{>{\raggedright\arraybackslash}p{3.2cm} >{\raggedright\arraybackslash}X >{\raggedright\arraybackslash}p{2.8cm}}
\toprule
Modul & Tujuan Pembelajaran & Output \\
\midrule
Cipher Klasik & Memahami transformasi teks asli ke teks sandi sederhana & Enkripsi dan dekripsi dasar \\
Diffie-Hellman & Memahami pertukaran kunci di kanal publik & Rahasia bersama \\
Pohon Merkle & Verifikasi integritas data berbasis hash & Hash akar dan bukti keanggotaan \\
\bottomrule
\end{tabularx}
\end{table}

Protokol Diffie-Hellman menggunakan eksponensiasi modular. Misalkan dua pihak setuju pada bilangan prima \(p\) dan generator \(g\). Masing-masing memilih rahasia \(a\) dan \(b\), lalu mengirim:
\[
A = g^a \mod p, \quad B = g^b \mod p
\]
Rahasia bersama dihitung sebagai:
\[
K = B^a \mod p = A^b \mod p = g^{ab} \mod p
\]

Sedangkan pohon Merkle menggunakan fungsi hash \(H(\cdot)\). Hash akar dihitung dengan menggabungkan hash anak-anak:
\[
H_{\text{root}} = H\big( H_{\text{kiri}} \parallel H_{\text{kanan}} \big)
\]

\section{Bagan Alir Pembelajaran Terstruktur}
\noindent
Diagram berikut menggambarkan siklus belajar adaptif: peserta dapat mengulang modul ketika
indikator pemahaman belum tercapai, sehingga proses belajar tetap bertahap dan terukur.

\begin{figure}[H]
\centering
\begin{tikzpicture}[
  node distance=1.6cm,
  block/.style={rectangle, rounded corners, draw=black, fill=blue!8, minimum width=5cm, minimum height=0.9cm, align=center},
  decision/.style={diamond, draw=black, fill=orange!12, aspect=2, align=center, inner sep=1pt},
  line/.style={-{Latex[length=2mm]}, thick}
]
\node[block] (start) {Mulai Belajar ALGORUST};
\node[block, below=of start] (fund) {Pelajari Algoritma Fundamental};
\node[decision, below=of fund] (check1) {Sudah paham\\kompleksitas?};
\node[block, below=of check1] (p2p) {Lanjut ke Modul P2P};
\node[block, below=of p2p] (crypto) {Lanjut ke Modul Kriptografi};
\node[decision, below=of crypto] (check2) {Semua latihan\\lulus?};
\node[block, below=of check2] (end) {Selesai dan Tinjau Portofolio Solusi};
\node[block, right=2.8cm of check1] (repeat1) {Ulangi materi\\dan latihan};
\node[block, right=2.8cm of check2] (repeat2) {Perbaiki solusi\\latihan};

\draw[line] (start) -- (fund);
\draw[line] (fund) -- (check1);
\draw[line] (check1) -- node[left]{Ya} (p2p);
\draw[line] (p2p) -- (crypto);
\draw[line] (crypto) -- (check2);
\draw[line] (check2) -- node[left]{Ya} (end);
\draw[line] (check1.east) -- node[above]{Tidak} (repeat1.west);
\draw[line] (repeat1.south) |- (fund.east);
\draw[line] (check2.east) -- node[above]{Tidak} (repeat2.west);
\draw[line] (repeat2.north) |- (crypto.east);
\end{tikzpicture}
\caption{Bagan alir jalur pembelajaran ALGORUST}
\end{figure}

\section{Contoh Implementasi Rust}
Contoh berikut menunjukkan implementasi BFS yang ringkas dan mudah diikuti, dengan format
penyorotan sintaks khusus Rust.

\begin{lstlisting}[style=ruststyle,caption={Contoh BFS sederhana di Rust}]
use std::collections::VecDeque;

fn bfs(graph: &[Vec<usize>], start: usize) -> Vec<usize> {
    let mut visited = vec![false; graph.len()];
    let mut order = Vec::new();
    let mut queue = VecDeque::new();

    visited[start] = true;
    queue.push_back(start);

    while let Some(node) = queue.pop_front() {
        order.push(node);
        for &next in &graph[node] {
            if !visited[next] {
                visited[next] = true;
                queue.push_back(next);
            }
        }
    }

    order
}
\end{lstlisting}

Berikut adalah contoh implementasi Caesar cipher di Rust untuk mengenkripsi teks:

\begin{lstlisting}[style=ruststyle,caption={Caesar cipher di Rust}]
fn caesar_encrypt(plain: &str, shift: u8) -> String {
    plain.chars()
        .map(|c| {
            if c.is_ascii_alphabetic() {
                let first = if c.is_ascii_lowercase() { b'a' } else { b'A' };
                let offset = (c as u8 - first + shift) % 26;
                (first + offset) as char
            } else {
                c
            }
        })
        .collect()
}
\end{lstlisting}

\section{Metodologi Evaluasi}
Pendekatan evaluasi dirancang untuk memastikan pemahaman konsep sekaligus kualitas implementasi:
\begin{itemize}
  \item Uji unit (\textit{unit test}) untuk validasi ketepatan keluaran.
  \item Latihan bertahap berbasis tingkat kesulitan.
  \item Analisis kompleksitas waktu dan ruang pada setiap solusi.
  \item Tinjauan kualitas kode: keterbacaan, modularitas, dan penanganan kasus tepi.
\end{itemize}

\section{Prinsip Keamanan dan Batasan}
Implementasi kriptografi dalam ALGORUST berorientasi edukasi, bukan produksi.
Untuk penggunaan nyata, disarankan:
\begin{itemize}
  \item Menggunakan pustaka kriptografi standar yang telah diaudit.
  \item Menghindari merancang algoritma kriptografi sendiri tanpa evaluasi pakar.
  \item Menerapkan pengelolaan kunci dan audit keamanan berkala.
\end{itemize}

\section{Penutup}
ALGORUST menempatkan pembelajaran algoritma dalam konteks rekayasa perangkat lunak modern:
berbasis data, terukur, dan berulang. Dengan pendekatan ini, peserta tidak hanya memahami
``cara kerja'' algoritma, tetapi juga ``kapan'' dan ``mengapa'' sebuah solusi dipilih.

\end{document}